Let \(R=(Y(0),Y(1))\).  We first verify the entropy range.  By the definition of \(h\) and the definition of \(\kappa^\dagger\),
\[
10\kappa^\dagger
=9\log\frac{10}{9}+\log 20
=\log\!\left(20\left(\frac{10}{9}\right)^9\right).
\]
Both terms defining \(\kappa^\dagger\) are nonnegative and \(h(9/10)>0\), so \(\kappa^\dagger>0\).  Moreover,
\[
20\cdot 10^9=20{,}000{,}000{,}000
<396{,}718{,}580{,}736=2^{10}9^9.
\]
Strict monotonicity of the logarithm therefore gives \(10\kappa^\dagger<10\log 2\), and hence
\[
0<\kappa^\dagger<\log 2.
\tag{1}
\]

We next give a constructive feasible three-state factorization.  Label the states \(00,10,11\), and set
\[
(\pi_{00},\pi_{10},\pi_{11})
=\left(\frac1{20},\frac9{10},\frac1{20}\right),
\qquad
(e_{00},e_{10},e_{11})=\left(0,\frac12,1\right).
\tag{2}
\]
At state \(ij\), let \(r_{ij\mid ij}=1\) and let every other response-type probability be zero.  Thus each \(r_{\cdot\cdot\mid u}\) belongs to \(\Delta_4\), \(\pi\in\Delta_3\), and \(e\in[0,1]^3\).  Substitution into the factual constraints in \(\mathcal F_{3,\kappa^\dagger}(P^\dagger)\) gives
\[
\begin{aligned}
\sum_u\pi_u(1-e_u)\sum_{i,j}r_{ij\mid u}\mathbf 1\{i=0\}
&=\frac1{20}=p^\dagger_{00},\\
\sum_u\pi_u(1-e_u)\sum_{i,j}r_{ij\mid u}\mathbf 1\{i=1\}
&=\frac9{20}=p^\dagger_{01},\\
\sum_u\pi_u e_u\sum_{i,j}r_{ij\mid u}\mathbf 1\{j=0\}
&=\frac9{20}=p^\dagger_{10},\\
\sum_u\pi_u e_u\sum_{i,j}r_{ij\mid u}\mathbf 1\{j=1\}
&=\frac1{20}=p^\dagger_{11}.
\end{aligned}
\tag{3}
\]
The entropy of the masses in (2) is
\[
-2\frac1{20}\log\frac1{20}-\frac9{10}\log\frac9{10}
=h\!\left(\frac9{10}\right)+\frac1{10}\log2
=\kappa^\dagger.
\tag{4}
\]
Finally, the definition of \(\Phi_3\) and (2) give
\[
\Phi_3(\pi,e,r)
=\left(\sum_u\pi_u r_{01\mid u},\sum_u\pi_u r_{10\mid u}\right)
=\left(0,\frac9{10}\right)
=\theta^\dagger.
\tag{5}
\]
Equations (2)--(5) prove
\[
\theta^\dagger\in\Phi_3\bigl(\mathcal F_{3,\kappa^\dagger}(P^\dagger)\bigr).
\tag{6}
\]
Equivalently, the constructive direction of \(\texttt{lem:factor-program-equivalence}\) turns this tuple into a factual-compatible law satisfying consistency, latent exchangeability, and the entropy budget.

It remains to rule out two positive latent states.  Consider any factual-compatible law having \((B,H)=\theta^\dagger\), and write \(J\) for its joint law of \((A,R)\).  Since response-type probabilities are nonnegative, the definitions of \(B\) and \(H\) imply
\[
J(R=01)=0,
\qquad
J(R=10)=\frac9{10},
\qquad
J(R=00)+J(R=11)=\frac1{10}.
\tag{7}
\]
By consistency, when \(A=0\) the observed outcome is the first component of \(R\), and when \(A=1\) it is the second component.  Factual compatibility with \(P^\dagger\), together with \(J(R=01)=0\), therefore yields
\[
J(A=0,R=00)=p^\dagger_{00}=\frac1{20},
\qquad
J(A=1,R=11)=p^\dagger_{11}=\frac1{20}.
\tag{8}
\]
The two displayed masses exhaust the total diagonal-type mass \(1/10\) in (7).  Nonnegativity consequently forces
\[
J(A=1,R=00)=J(A=0,R=11)=0.
\tag{9}
\]
Applying factual compatibility to the other two observed cells and using (9) gives
\[
J(A=0,R=10)=p^\dagger_{01}=\frac9{20},
\qquad
J(A=1,R=10)=p^\dagger_{10}=\frac9{20}.
\tag{10}
\]
Thus (7)--(10) force the complete joint table
\[
J(0,00)=\frac1{20},\quad
J(0,10)=\frac9{20},\quad
J(1,10)=\frac9{20},\quad
J(1,11)=\frac1{20},
\tag{11}
\]
with every other cell equal to zero.

Suppose for contradiction that a tuple in \(\mathcal F_{2,\kappa^\dagger}(P^\dagger)\) has image \(\theta^\dagger\).  By the constructive equivalence in \(\texttt{lem:factor-program-equivalence}\), its positive-mass latent states give a product-mixture representation
\[
J(a,r)=\sum_w\lambda_w e_w^a(1-e_w)^{1-a}s_w(r),
\tag{12}
\]
where \(\lambda_w>0\) on the retained support, \(e_w\in[0,1]\), \(s_w\in\Delta_4\), and there are at most two retained states.  Positivity of \(J(0,00)\) in (11) provides a state \(w_0\) for which \(\lambda_{w_0}s_{w_0}(00)>0\) and \(1-e_{w_0}>0\).  But
\[
0=J(1,00)=\sum_w\lambda_w e_ws_w(00),
\]
and all summands are nonnegative, so \(e_{w_0}=0\).  Similarly, positivity of \(J(1,11)\) provides a state \(w_1\) with \(\lambda_{w_1}s_{w_1}(11)>0\), whereas
\[
0=J(0,11)=\sum_w\lambda_w(1-e_w)s_w(11)
\]
forces \(e_{w_1}=1\).  Hence \(w_0\ne w_1\).  Because (12) has at most two positive-mass states, these are all of them, and the latent state determines \(A\).  The marginal treatment law under \(P^\dagger\) is balanced, since
\[
P^\dagger(A=0)=p^\dagger_{00}+p^\dagger_{01}=\frac12,
\qquad
P^\dagger(A=1)=p^\dagger_{10}+p^\dagger_{11}=\frac12.
\]
It follows that \(\lambda_{w_0}=\lambda_{w_1}=1/2\), so the entropy of the factor in (12) is
\[
-\frac12\log\frac12-\frac12\log\frac12=\log2>\kappa^\dagger
\tag{13}
\]
by (1).  This contradicts the entropy constraint in \(\mathcal F_{2,\kappa^\dagger}(P^\dagger)\).  Therefore
\[
\theta^\dagger\notin\Phi_2\bigl(\mathcal F_{2,\kappa^\dagger}(P^\dagger)\bigr).
\tag{14}
\]
This zero-pattern product-mixture argument is the response-type embedding of the nonrepresentability technique in \citet[Section V, Example 2 and Appendix A]{KumarLiElGamal2014}; all properties needed here have been derived directly in (7)--(13).

Finally, any element \((\pi,e,r)\in\mathcal F_{2,\kappa^\dagger}(P^\dagger)\) can be padded by a third state with mass zero and arbitrary admissible propensity and response distribution.  The convention \(0\log0=0\) shows that padding preserves entropy, and the zero mass shows that it preserves every factual constraint and \(\Phi_m\).  Hence
\[
\Phi_2\bigl(\mathcal F_{2,\kappa^\dagger}(P^\dagger)\bigr)
\subseteq
\Phi_3\bigl(\mathcal F_{3,\kappa^\dagger}(P^\dagger)\bigr).
\]
Together with (6) and (14), this inclusion is strict, proving every assertion of the proposition.